% =============================================================================
% Baily-Borel cusp diagrams: constructors
% =============================================================================
% Building blocks for cusp diagrams of Baily-Borel compactifications. A
% constructor takes the data of cusps and incidences and draws them; the cusp
% diagrams themselves are objects, one file each, under figures/objects/bb-cusps/,
% assembled from these constructors. A 0-cusp (the orbit of a primitive
% isotropic line) is a cusp0 node and a 1-cusp (the orbit of a primitive
% isotropic plane) a cusp1 node; an incidence arrow runs from a 0-cusp to each
% 1-cusp whose closure contains it. A 0-cusp or 1-cusp is named by its boundary
% lattice I^perp/I; where the lattice has a standard name it is used (E_{10}(2)
% = U(2)+E_8(2), D_8^*(2), D_{16}^+, A_1^{+7}). Place a diagram with
%   \pic (S) {object=bb-cusps/fen2};
% and decorate through the node centres, e.g. the image of another boundary:
%   \scoped[on background layer] \draw[cusp image] (S-2.center) -- (S-245.center);
% \input by dzg-constructors.tex, so "@" is a letter here.
%
% Commands:
%   \bbcusp{<cusp0|cusp1, options>}{<name>}{<coordinate>}{<content>}
%          {<label direction>}{<canonical label>}{<eta label>}
%       a cusp node with content $<content>$ and the label of the active
%       scheme, $<canonical label>$ or $<eta label>$; an empty label in the
%       active scheme draws none.
%   \bblabel{<canonical>}{<eta>}   the text of the active label scheme (empty
%                                  under cusp labels=none), for node contents
%   \bbincidences{<0-cusp>/<1-cusp>/<divisibility>, ...}
%       incidence arrows; divisibility is that of the second isotropic vector
%       of the plane in the 0-cusp lattice, 1 or 2 (an even incidence).
%   \bbbrace{<x from>}{<x to>}{<y>}{<text>}
%       a brace under a column with the invariants $<text>$ of its lattices,
%       drawn in the canonical scheme only.
%   \bbcolumn{<x>}{<name prefix>}{<first>}{<second>}{<next to last>}{<last>}
%       a column of four 1-cusps at x, y = 2, 1, -1, -2, named <prefix>1, ...,
%       <prefix>4, with a vertical ellipsis at y = 0: a family C_1, ..., C_n.
%
% Keys:
%   cusp labels=canonical   the canonical labels of each object (default)
%   cusp labels=eta         the names eta_i, I_{i,j} of the dissertation text
%                           (Enriques-K3 chapter) where the object has them
%   cusp labels=none        no labels
%   incidence divisibility  double each incidence whose second isotropic
%                           vector has divisibility 2 in the 0-cusp lattice
%                           (an even move); a single arrow is divisibility 1.
%                           For 2-elementary lattices this is the incidence at
%                           which the a-invariant drops by 2 [AE22, Prop. 5.5].

\def\dzg@bb@scheme{canonical}
\tikzset{
  cusp labels/.is choice,
  cusp labels/canonical/.code={\def\dzg@bb@scheme{canonical}},
  cusp labels/eta/.code={\def\dzg@bb@scheme{eta}},
  cusp labels/none/.code={\def\dzg@bb@scheme{none}},
  incidence divisibility/.style={even incidence/.style={doubled mark}},
  even incidence/.style={},
}

\def\dzg@bb@canonical{canonical}
\def\dzg@bb@eta{eta}
\def\dzg@bb@none{none}
\newcommand\bblabel[2]{\csname dzg@bb@label@\dzg@bb@scheme\endcsname{#1}{#2}}
\def\dzg@bb@label@canonical#1#2{#1}
\def\dzg@bb@label@eta#1#2{#2}
\def\dzg@bb@label@none#1#2{}

\newcommand\bbcusp[7]{%
  \ifx\dzg@bb@scheme\dzg@bb@canonical\def\dzg@bb@text{#6}\fi
  \ifx\dzg@bb@scheme\dzg@bb@eta\def\dzg@bb@text{#7}\fi
  \ifx\dzg@bb@scheme\dzg@bb@none\let\dzg@bb@text\@empty\fi
  \ifx\dzg@bb@text\@empty
    \node[#1] (#2) at (#3) {$#4$};%
  \else
    \node[#1, label={[cusp label]#5:$\dzg@bb@text$}] (#2) at (#3) {$#4$};%
  \fi}

\newcommand\bbincidences[1]{%
  \foreach \dzg@p/\dzg@c/\dzg@d in {#1}{%
    \ifnum\dzg@d=2
      \draw[incidence, even incidence] (\dzg@p) -- (\dzg@c);%
    \else
      \draw[incidence] (\dzg@p) -- (\dzg@c);%
    \fi}}

\newcommand\bbbrace[4]{%
  \ifx\dzg@bb@scheme\dzg@bb@canonical
    \draw[cusp brace] (#1,#3) -- (#2,#3) node[midway, below=5.5ex] {$#4$};%
  \fi}

\newcommand\bbcolumn[6]{%
  \foreach \dzg@k/\dzg@y/\dzg@L in {1/2/{#3}, 2/1/{#4}, 3/-1/{#5}, 4/-2/{#6}}
    {\bbcusp{cusp1}{#2\dzg@k}{#1,\dzg@y}{\dzg@L}{}{}{}}%
  \node at (#1,0) {$\vdots$};}
